\documentclass[10pt,letterpaper]{article}

% Preprint layout: full-width title and abstract, then a two-column body.
% Not a conference submission and not an arXiv deposit.
\usepackage[T1]{fontenc}
\usepackage{mathptmx}
\usepackage{amsmath}
\usepackage[letterpaper,margin=0.75in,columnsep=0.3in,headheight=14.5pt,headsep=10pt,footskip=28pt]{geometry}
\usepackage{fancyhdr}
\usepackage[font=small,labelfont=bf,labelsep=colon]{caption}
\usepackage{microtype}

\pagestyle{fancy}
\fancyhf{}
\fancyhead[L]{\footnotesize\itshape A Per-Gaussian Physics-Prior MLP}
\fancyhead[R]{\footnotesize Preprint}
\fancyfoot[C]{\thepage}
\renewcommand{\headrulewidth}{0.4pt}
\fancypagestyle{plain}{%
  \fancyhf{}%
  \fancyfoot[L]{\footnotesize Preprint}%
  \fancyfoot[C]{\thepage}%
  \renewcommand{\headrulewidth}{0pt}%
}

\makeatletter
\renewcommand{\@seccntformat}[1]{\csname the#1\endcsname.\enspace}
\renewcommand\section{\@startsection{section}{1}{\z@}%
  {-2.4ex \@plus -0.6ex \@minus -0.2ex}%
  {1.1ex \@plus 0.2ex}%
  {\large\bfseries\raggedright}}
\let\savedthebibliography\thebibliography
\renewcommand{\thebibliography}[1]{%
  \savedthebibliography{#1}%
  \small
}
\makeatother

\DeclareRobustCommand{\_}{\textunderscore\allowbreak}
\frenchspacing
\emergencystretch=2em
\sloppy
\newsavebox\titlebox

\begin{document}

\begin{lrbox}{\titlebox}
\begin{minipage}{\textwidth}
\begin{center}
{\LARGE\bfseries A Per-Gaussian Physics-Prior MLP\\[0.15em]
Trained on Fitted Splats\par}
\vspace{0.7em}
{\large sandy-t-shelvan\par}
\vspace{0.35em}
{9 November 2026\par}
\vspace{0.25em}
{\footnotesize Preprint\par}
\end{center}
\vspace{0.15em}
\begin{center}
{\large\bfseries Abstract}
\end{center}
\begin{quote}
A per-Gaussian multilayer perceptron predicts physical attributes of an object from a fitted 3D Gaussian splat~\cite{kerbl2023gaussian}.
The network has two tanh hidden layers of width 32 and a linear head of width 33.
Inference runs the MLP on each in-mask Gaussian and averages the 33 outputs.
It does not average the 912 inputs.
The six weight arrays in \texttt{data/mlp/physics\_prior\_splat.npz} were trained for this project.
They are not a released third-party checkpoint.
The saved weights are epoch 10 of a 14-epoch run on 3274 objects.
On the 323-object test split, which was scored once after the checkpoint was chosen, material accuracy is 0.8885448916408669, body accuracy is 0.9907120743034056, and collision accuracy is 0.7678018575851393.
Meshes and per-Gaussian \texttt{prior\_input.npy} files are not included.
\end{quote}
\vspace{0.35em}
\end{minipage}
\end{lrbox}

\twocolumn[\usebox\titlebox]
\thispagestyle{plain}

\section{Task}
The model reads one fitted splat and predicts a material class, a body type, a collision primitive, a log mass, a log density, a rigidity, and a deformability.
Each Gaussian is one row.
The MLP is applied to every in-mask Gaussian, and the 33 raw outputs are averaged.
Averaging the 912-dimensional inputs and running the MLP once is a different function and is not the inference used here.

A material vote that ends in a tie returns no class.
Keyword proposals are passed in by the caller.
The vote does not look up the keyword dictionary on its own.
The unknown class weight is multiplied by 0.4 before it is added to the class score.
An object whose deformability is at least 0.65 is resolved as a soft body.
The published body accuracy in Section~\ref{sec:results} is the body head compared with the stored label.
The soft-body rule is applied when a physics record is built, not inside the masked training loss.

The weights were trained in this project.
They correspond to the per-Gaussian physics prior in Section 3.4 of Hollis, Fan, and Yang~\cite{hollis2026sceneagent}.
They are not a released third-party checkpoint.
The keyword dictionary is the table in \texttt{sceneagent/taxonomy.py} (\texttt{\_KEYWORD\_RULES}).
It is this project's dictionary.

\section{Input}
One row has width 912.
Columns 0--383 are DINOv2 ViT-S/14~\cite{oquab2024dinov2}, width 384 (\texttt{vit\_small\_patch14\_dinov2.lvd142m}).
Columns 384--895 are OpenCLIP ViT-B-16~\cite{cherti2023reproducible} with the OpenAI CLIP weights~\cite{radford2021clip}, width 512.
Columns 896--911 are sixteen geometry quantities, in this order: \(x\), \(y\), \(z\), three exp scales, sigmoid opacity, quaternion components \texttt{rot\_0} through \texttt{rot\_3}, spherical-harmonic DC coefficients \texttt{f\_dc\_0} through \texttt{f\_dc\_2}, log volume, and anisotropy.
Log volume is the sum of the three log-scales stored in the ply.
Anisotropy is the ratio of the largest exp scale to the smallest exp scale.

A ply that lacks a required source column is refused.
A Gaussian whose source values are not finite is dropped.
Missing numbers are not filled.
Opacity must vary.
A fitted row is kept only when the opacity standard deviation is at least \(10^{-3}\) and the opacity is not the constant 0.99.
A triangle prior with opacity fixed at 0.99 is the wrong input.
That stand-in is not this model, and the file \texttt{physics\_prior.npz} is not part of this repository.

DINOv2 and OpenCLIP stay frozen.
They are not parameters of Adam~\cite{kingma2015adam}.
splatfacto, the Nerfstudio method used here~\cite{tancik2023nerfstudio}, fits the Gaussians and is not updated by Adam.

\section{Output}
The head has width 33.
\begin{itemize}
\item Slots 0--9 are material logits: plastic, metal, wood, fabric, paper and cardboard, rubber and foam, glass and ceramic, food and organic matter, liquid and granular matter, unknown.
\item Slots 10--15 are body logits: dynamic rigid, static rigid, kinematic, soft deformable, cloth shell, fluid or granular.
\item Slots 16--23 are collision logits: box, sphere, cylinder, capsule, convex hull, triangle mesh, tetrahedral mesh, cloth mesh.
\item Slot 24 is log mass and slot 25 is log density, both natural logs.
\item Slot 26 is rigidity and slot 27 is deformability.
The sigmoid is not inside the forward pass.
It is applied to these two slots in the loss and when the mean is decoded.
\item Slots 28--32 are numeric fields that stay masked.
They receive no target and no gradient.
\end{itemize}
Friction and restitution are null.
Young's modulus, Poisson's ratio, and the rendering material are unpublished and were not invented.

\section{Network}
The network is a multilayer perceptron.
It is not a recurrent network and not a convolutional network.
Two hidden layers, each of width 32, use \(\tanh\).
The head is linear.
The six trained arrays, with Glorot uniform initialization~\cite{glorot2010understanding} and zero biases, are
\begin{center}
\begin{tabular}{ll}
\(w_1\) & \((912, 32)\) \\
\(b_1\) & \((32,)\) \\
\(w_2\) & \((32, 32)\) \\
\(b_2\) & \((32,)\) \\
\(w_{\mathrm{out}}\) & \((32, 33)\) \\
\(b_{\mathrm{out}}\) & \((33,)\) \\
\end{tabular}
\end{center}
Glorot uniform draws each entry from \([-\ell, \ell]\) with \(\ell = \sqrt{6/(n_{\mathrm{in}}+n_{\mathrm{out}})}\).
Only these six arrays are trained.
The forward map is
\begin{align*}
h_1 &= \tanh(x W_1 + b_1), \\
h_2 &= \tanh(h_1 W_2 + b_2), \\
y &= h_2 W_{\mathrm{out}} + b_{\mathrm{out}}.
\end{align*}
Rows are accumulated in chunks of 4096 Gaussians.
The chunked mean matches the mean of the full forward.

\section{Loss}
The loss is applied to the in-mask mean, one object at a time.
Material, body, and collision use cross-entropy on the logits.
Log mass and log density use squared error on the raw outputs.
Rigidity and deformability use squared error after the sigmoid.
Slots 28--32 are masked, so their contribution is zero.
For one object the loss is
\begin{align*}
\mathcal{L}
&= \mathrm{CE}(y_{0:9}, \mathrm{material})
+ \mathrm{CE}(y_{10:15}, \mathrm{body}) \\
&\quad + \mathrm{CE}(y_{16:23}, \mathrm{collision})
+ (y_{24} - \log m)^2 \\
&\quad + (y_{25} - \log \rho)^2
+ (\sigma(y_{26}) - r)^2
+ (\sigma(y_{27}) - d)^2.
\end{align*}
Class accuracy is the fraction of objects whose argmax matches the label.
Numeric error is the mean absolute error of the decoded log mass, log density, sigmoid rigidity, and sigmoid deformability.

\section{Optimizer}
The optimizer is a hand-written Adam.
The learning rate is \(10^{-3}\), \(\beta_1 = 0.9\), \(\beta_2 = 0.999\), and \(\epsilon = 10^{-8}\).
Weight decay of \(10^{-6}\) is added to the gradient before the moments are updated.
A global \(\ell_2\) clip of 5.0 is applied to the six gradients before that decay term.
One object is one step.
The seed is 0.
The maximum is 30 epochs.
From epoch 12, training stops when the best of the last 4 validation losses is not within 0.5 percent of the best loss before that window.
The threshold is 0.995, and the comparison is strict, so a drop of exactly 0.5 percent does not stop.
The saved weights are the earliest epoch with the lowest validation loss.
The test split is not read by the stopping rule.

\section{Data}
Each fitted object is built on the CPU and then packed.
The orbit is 16 views: 8 azimuths at elevations \(18^\circ\) and \(50^\circ\).
COLMAP~\cite{schoenberger2016sfm} registers those views, splatfacto runs for 600 steps, and the pack writes \((N, 912)\).
Labels come from \texttt{gemini-3.8-flash}, with thinking level \texttt{low} and temperature 0.
The labeler caps itself at 500 requests per minute, 1{,}500{,}000 tokens per minute, and 8{,}000 requests per day.
Ids that already had an ok label were not relabeled.

The queue is about 7538 objects: YCB~\cite{calli2017ycb} 71, GSO~\cite{downs2022gso} 1024, and ABO~\cite{collins2022abo} 6443.
Ok labels used as the original targets number 4427: YCB 70, GSO 244, and ABO 4113.
Those rows are \texttt{data/mlp/stratified\_labels.4090.jsonl}.
The fitted rows in this training run number 3736: YCB 0, GSO 162, and ABO 3574.
About 191 fits failed, mostly because COLMAP wrote no model.
A later file, \texttt{data/mlp/stratified\_labels.remaining.jsonl}, has 1796 additional ok labels (GSO 737, ABO 1059).
That file was not used to train this checkpoint.

\section{Split}
Train ids are the complement of a 12 percent identity holdout.
Three known ids are held out when they are present: \texttt{002\_master\_chef\_can}, \texttt{Weisshai\_Great\_White\_Shark}, and \texttt{B079TXJCR5}.
Those three were absent from the fitted table, so the holdout is the hash sample alone.
An id is selected when the first byte of \(\mathrm{sha256}(\mathrm{id})\), divided by 256, is below 0.12.
The 462 holdout ids are sorted by the sha256 digest and then by the id string.
Validation is the first 30 percent, half up, using integer arithmetic: \(\mathrm{divmod}(3n, 10)\), with a remainder of 5 or more rounding up.
For \(n = 462\) the quotient is 138 and the remainder is 6, so validation has 139 ids and test has 323.
Train has 3274.
Seed 0 does not shuffle this validation slice.
It does permute the order of training objects inside each epoch.
The test ids are scored once, after the checkpoint is chosen.

\section{Results}
\label{sec:results}
The report is \texttt{data/mlp/train\_report\_splat.json} (375162 bytes).
The checkpoint is \texttt{data/mlp/physics\_prior\_splat.npz} (260290 bytes).
The run recorded \texttt{epochs\_ran} 14, \texttt{checkpoint\_epoch} 10, and \texttt{n\_fitted} 3736.
Table~\ref{tab:scores} is \texttt{score\_split} on the restored epoch-10 weights.
The in-loop losses in Table~\ref{tab:curve} are the mean training loss while the weights were still moving, and the validation loss of the weights at the end of that epoch.
They are not the restored-checkpoint scores.

\begin{table*}[t]
\centering
\caption{Restored epoch-10 checkpoint. Classification entries are accuracy. Numeric entries are mean absolute error.}
\label{tab:scores}
\small
\setlength{\tabcolsep}{5pt}
\begin{tabular}{lrrr}
\hline
 & Train & Validation & Test \\
\hline
\(n\) & 3274 & 139 & 323 \\
Loss & 1.3504881230964794 & 2.293009738617003 & 2.0879311734829393 \\
Material & 0.9526572999389127 & 0.9280575539568345 & 0.8885448916408669 \\
Body & 0.9932803909590715 & 0.9568345323741008 & 0.9907120743034056 \\
Collision & 0.8313989004276114 & 0.8201438848920863 & 0.7678018575851393 \\
Log mass & 0.29221152206794304 & 0.38314453113056723 & 0.35213020169288767 \\
Log density & 0.493749474065539 & 0.5773056012440149 & 0.5719236185569908 \\
Rigidity & 0.03465716247218895 & 0.05072960603976827 & 0.037416784396735514 \\
Deformability & 0.03273072179495256 & 0.046840128143126915 & 0.03501937974941083 \\
\hline
\end{tabular}
\end{table*}

\begin{table*}[t]
\centering
\caption{In-loop mean loss by epoch on the 3274 training objects, and validation loss on the 139 validation objects. Epoch 10 has the lowest validation loss.}
\label{tab:curve}
\small
\setlength{\tabcolsep}{8pt}
\begin{tabular}{rrr}
\hline
Epoch & Train loss & Validation loss \\
\hline
1 & 3.459843983024622 & 2.908307250704289 \\
2 & 2.2695306359359195 & 2.8169932549644834 \\
3 & 2.033521992464412 & 2.652548772876703 \\
4 & 1.8945112152633685 & 2.883813949546511 \\
5 & 1.8108405343027458 & 2.5052527957759665 \\
6 & 1.6972543350052578 & 2.353592547470898 \\
7 & 1.6250525286905761 & 2.58225741753844 \\
8 & 1.5474230001040072 & 2.521480992803893 \\
9 & 1.4982185734968427 & 2.3677041720399106 \\
10 & 1.4523404458332272 & 2.293009738617003 \\
11 & 1.3954204487159254 & 2.4070450118198785 \\
12 & 1.3562194154564946 & 2.4512925276327464 \\
13 & 1.3246501546116667 & 2.4344396546699816 \\
14 & 1.2731846908738063 & 2.386788401613216 \\
\hline
\end{tabular}
\end{table*}

In-loop train loss fell from 3.459843983024622 at epoch 1 to 1.2731846908738063 at epoch 14.
Epoch 10 of that curve is 1.4523404458332272.
Validation loss at epoch 14 was 2.386788401613216, higher than 2.293009738617003 at epoch 10, so the earlier epoch was kept.
More epochs on this same 3274-object set did not improve validation.
Early stopping ended the run at epoch 14.
The best validation loss in the report is 2.293009738617003.

\section{Previous run}
An earlier 30-epoch run trained on the same 3274 fitted objects and scored the whole 462-id holdout, with no validation split.
It is a previous run, not the result of this paper.
The backup files are \texttt{data/mlp/physics\_prior\_splat\_holdout462.npz} and \texttt{data/mlp/train\_report\_splat\_holdout462.json}.
That report records holdout loss 2.3994417890371444, material accuracy 0.8939393939393939 (0.894 to three digits), body accuracy 0.9783549783549783 (0.978), and collision accuracy 0.7770562770562771 (0.777).
The holdout mean absolute errors are log mass 0.36483208720055876 (0.365), log density 0.5572938724288687 (0.557), rigidity 0.039427355866847605 (0.039), and deformability 0.03778218721477836 (0.038).
Train material accuracy on that run is 0.9813683567501528 (0.981).

\section{Limitations}
Some ABO masses in the 4427-row label file are extreme.
Sixty-six ABO rows have \texttt{mass\_kg} above 100, and the largest value in that file is 580.
Fifty-six rows are named the placeholder \texttt{consolidatedScene}.
Forty-three of the 3736 fitted ids have that name.
Cloth-shell rows with deformability at least 0.65 would be rewritten to soft deformable by the body resolver.
The 4427-row file contains 535 such rows; the stored labels were left as written.
About 1315 queue ids still had no ok label when the 1796-row file was written, mostly because the mesh was the wrong size or the render failed.
The 1796 later labels were not used to train the epoch-10 checkpoint.
Rescoring the checkpoint requires the per-Gaussian \texttt{prior\_input.npy} arrays, which are not in this repository.

\begin{thebibliography}{99}

\bibitem{kerbl2023gaussian}
Bernhard Kerbl, Georgios Kopanas, Thomas Leimk{\"u}hler, and George Drettakis.
3D Gaussian Splatting for Real-Time Radiance Field Rendering.
\textit{ACM Transactions on Graphics}, 42(4), 2023.
Proceedings of SIGGRAPH.

\bibitem{schoenberger2016sfm}
Johannes~L. Sch{\"o}nberger and Jan-Michael Frahm.
Structure-from-Motion Revisited.
In \textit{Proceedings of the IEEE Conference on Computer Vision and Pattern Recognition (CVPR)}, pages 4104--4113, 2016.

\bibitem{tancik2023nerfstudio}
Matthew Tancik, Ethan Weber, Evonne Ng, Ruilong Li, Brent Yi, Justin Kerr, Terrance Wang, Alexander Kristoffersen, Jake Austin, Kamyar Salahi, Abhik Ahuja, David McAllister, and Angjoo Kanazawa.
Nerfstudio: A Modular Framework for Neural Radiance Field Development.
In \textit{ACM SIGGRAPH 2023 Conference Proceedings}, 2023.

\bibitem{oquab2024dinov2}
Maxime Oquab, Timoth{\'e}e Darcet, Th{\'e}o Moutakanni, Huy~V. Vo, Marc Szafraniec, Vasil Khalidov, Pierre Fernandez, Daniel Haziza, Francisco Massa, Alaaeldin El-Nouby, Mahmoud Assran, Nicolas Ballas, Wojciech Galuba, Russell Howes, Po-Yao Huang, Shang-Wen Li, Ishan Misra, Michael Rabbat, Vasu Sharma, Gabriel Synnaeve, Hu Xu, Herv{\'e} J{\'e}gou, Julien Mairal, Patrick Labatut, Armand Joulin, and Piotr Bojanowski.
DINOv2: Learning Robust Visual Features without Supervision.
\textit{Transactions on Machine Learning Research}, 2024.
arXiv:2304.07193.

\bibitem{radford2021clip}
Alec Radford, Jong~Wook Kim, Chris Hallacy, Aditya Ramesh, Gabriel Goh, Sandhini Agarwal, Girish Sastry, Amanda Askell, Pamela Mishkin, Jack Clark, Gretchen Krueger, and Ilya Sutskever.
Learning Transferable Visual Models From Natural Language Supervision.
In \textit{Proceedings of the 38th International Conference on Machine Learning}, volume 139 of \textit{Proceedings of Machine Learning Research}, pages 8748--8763, 2021.

\bibitem{cherti2023reproducible}
Mehdi Cherti, Romain Beaumont, Ross Wightman, Mitchell Wortsman, Gabriel Ilharco, Cade Gordon, Christoph Schuhmann, Ludwig Schmidt, and Jenia Jitsev.
Reproducible Scaling Laws for Contrastive Language-Image Learning.
In \textit{Proceedings of the IEEE/CVF Conference on Computer Vision and Pattern Recognition (CVPR)}, pages 2818--2829, 2023.

\bibitem{kingma2015adam}
Diederik~P. Kingma and Jimmy Ba.
Adam: A Method for Stochastic Optimization.
In \textit{International Conference on Learning Representations}, 2015.

\bibitem{glorot2010understanding}
Xavier Glorot and Yoshua Bengio.
Understanding the difficulty of training deep feedforward neural networks.
In \textit{Proceedings of the Thirteenth International Conference on Artificial Intelligence and Statistics}, volume 9 of \textit{Proceedings of Machine Learning Research}, pages 249--256, 2010.

\bibitem{calli2017ycb}
Berk Calli, Arjun Singh, James Bruce, Aaron Walsman, Kurt Konolige, Siddhartha Srinivasa, Pieter Abbeel, and Aaron~M. Dollar.
Yale-CMU-Berkeley dataset for robotic manipulation research.
\textit{The International Journal of Robotics Research}, 36(3):261--268, 2017.

\bibitem{downs2022gso}
Laura Downs, Anthony Francis, Nate Koenig, Brandon Kinman, Ryan Hickman, Krista Reymann, Thomas~B. McHugh, and Vincent Vanhoucke.
Google Scanned Objects: A High-Quality Dataset of 3D Scanned Household Items.
In \textit{2022 International Conference on Robotics and Automation (ICRA)}, pages 2553--2560, 2022.

\bibitem{collins2022abo}
Jasmine Collins, Shubham Goel, Kenan Deng, Achleshwar Luthra, Leon Xu, Erhan Gundogdu, Xi Zhang, Tomas~F. Yago Vicente, Thomas Dideriksen, Himanshu Arora, Matthieu Guillaumin, and Jitendra Malik.
ABO: Dataset and Benchmarks for Real-World 3D Object Understanding.
In \textit{Proceedings of the IEEE/CVF Conference on Computer Vision and Pattern Recognition (CVPR)}, pages 21126--21136, 2022.

\bibitem{hollis2026sceneagent}
Hollis, Fan, and Yang.
SceneAgent: An Agentic Pipeline from 3D and 4D Captures to Simulation Environments.
Harvard Computational Robotics Group, preprint, 2026.
arXiv:2610.34745.

\end{thebibliography}

\end{document}
